\subsection{Convolution triples: construction and security}\label{sec:heformal}

This subsection defines the convolution triple, the packed encoding and the protocol that generates it, and argues
correctness, security and cost. The engineering steps that follow (\cref{fig:steps}) change how fast the
protocol runs, not what it sends.

\paragraph{Notation.} All arithmetic shares live in $\mathbb{Z}_{2^{32}}$. BFV works over
$R_t = \mathbb{Z}_t[X]/(X^N+1)$ with $N = 4096$ and plaintext modulus $t = 2^{32}$, so a plaintext coefficient
is exactly one ring element. The ciphertext modulus is $q = q_0 q_1$ with primes of 60 and 49 bits
(109 bits, 128-bit security by the HE standard's table), $\Delta = \lfloor q/t\rfloor$, $s$ is a party's
ternary secret key and $\chi$ the error distribution. $\mathrm{Enc}_s$ denotes secret-key and
$\mathrm{Enc}_{\mathit{pk}}$ public-key encryption. A convolution $y = W \star x$ is computed after the
polyphase reduction of \texttt{ConvLayout}: a stride-$\sigma$, padded convolution is rewritten as a stride-1,
unpadded convolution over $\sigma^2$ phase channels, which yields exactly the strided outputs. We therefore only
treat stride-1 valid convolutions of $B$ images of $C \times H \times W$ with $M$ filters of
$C \times k_h \times k_w$.

\paragraph{The triple, and why AB2 halves the traffic.} With masked values $m_v = v + \lambda_v$, the
ABY2 online step of a convolution is
\[
  m_y \;=\; m_W \star m_x \;-\; m_W \star \lambda_x \;-\; \lambda_W \star m_x \;+\; \Gamma \;+\; \lambda_y,
  \qquad \Gamma = \lambda_W \star \lambda_x ,
\]
computed share-wise (P0 adds $m_W \star m_x$). The preprocessing must deliver additive shares of $\Gamma$, the
\emph{convolution triple}. Write $\lambda_v = \lambda_v^0 + \lambda_v^1$ for the parties' shares.
\begin{itemize}\itemsep1pt
\item \textbf{AB} (\flag{A_KNOWN=0}). Both masks are shared, and
  $\Gamma = \lambda_W^0 \star \lambda_x^0 + \lambda_W^1 \star \lambda_x^1 + \lambda_W^0 \star \lambda_x^1 + \lambda_W^1 \star \lambda_x^0$.
  The first two terms are local; each cross term needs one secure product in which one party holds a filter and
  the other an input: two HE products per layer, one in each direction.
\item \textbf{AB2} (\flag{A_KNOWN=1}). The model owner P0 holds the whole weight mask ($\lambda_W^1 = 0$), and
  $\Gamma = \lambda_W \star (\lambda_x^0 + \lambda_x^1)$ is a \emph{single} product: P1 encrypts $\lambda_x^1$,
  P0 adds its own $\lambda_x^0$ inside the ciphertext and multiplies by its plaintext filters.
\end{itemize}
AB2 therefore needs half the ciphertexts, half the bytes and half the HE work of AB: 479 against 958\,MiB for
the ImageNet convolutions and 99 against 197\,MiB for 10 CIFAR-10 images (\cref{tab:concurrent}), with the same
protocol for each product.
It is secure because it is ABY2's own input semantics: the party that inputs a value chooses that value's whole
mask. The weights are P0's input. With \flag{SHARE_PREP} P0 sets $\lambda_W = -W$, so the public masked
weight is $m_W = 0$ and P1's share is 0: P1 learns nothing about $W$, and P0 learns only a mask of its own input.
Nothing about $x$ reaches P0, since $\lambda_x^1$ stays encrypted under P1's key. AB is needed only when no single
party owns the weights (e.g.\ a secret-shared model); with the weights known in preprocessing
(\flag{MODELWEIGHTS_KNOWN_DURING_PREPROCESSING}) the same single product is used, with P1's share of $\Gamma$
prescribed by its PRNG stream instead of drawn fresh, so no extra message is needed.

\paragraph{Packed encoding.} The encoding follows the block layout of troy's \texttt{Conv2dHelper}. A tiling is
$(b, c_i, c_o, h, w)$ with $b\,c_i\,c_o\,h\,w \le N$, $h \ge k_h$, $w \ge k_w$; a tile covers $b$ images and an
$h \times w$ window, and yields $y_h \times y_w = (h - k_h + 1) \times (w - k_w + 1)$ outputs per image and filter.
For a tile $\tau$ (first image $\beta_0$, first row $r_0$, first column $s_0$) and an input group $g$ of $c_i$
channels, and an output group $\gamma$ of $c_o$ filters,
\begin{align*}
  \pi_x^{\tau,g}(X) &= \sum_{\beta<b}\sum_{c<c_i}\sum_{i<h}\sum_{j<w}
     x[\beta_0{+}\beta,\; g c_i{+}c,\; r_0{+}i,\; s_0{+}j]\; X^{(\beta c_i c_o + c)\,hw + iw + j}, \\
  \pi_W^{\gamma,g}(X) &= \sum_{o<c_o}\sum_{c<c_i}\sum_{a<k_h}\sum_{d<k_w}
     W[\gamma c_o{+}o,\; g c_i{+}c,\; k_h{-}1{-}a,\; k_w{-}1{-}d]\; X^{(o c_i + c_i - 1 - c)\,hw + aw + d}.
\end{align*}
\textbf{Lemma 1.} For $\beta < b$, $o < c_o$, $i < y_h$, $j < y_w$ let
$e(\beta,o,i,j) = (\beta c_i c_o + o c_i + c_i - 1)\,hw + (k_h - 1 + i)\,w + (k_w - 1 + j)$. Then
$[\pi_x^{\tau,g}\cdot\pi_W^{\gamma,g}]_{e(\beta,o,i,j)} =
\sum_{c<c_i}\sum_{a<k_h}\sum_{d<k_w} x[\beta_0{+}\beta, gc_i{+}c, r_0{+}i{+}a, s_0{+}j{+}d]\cdot W[\gamma c_o{+}o, gc_i{+}c, a, d]$
in $R_t$, i.e.\ the input group's partial sum of output $(\beta_0{+}\beta,\gamma c_o{+}o, r_0{+}i, s_0{+}j)$, and
$\sum_g \pi_x^{\tau,g}\pi_W^{\gamma,g}$ holds the full output there.

\emph{Proof sketch.} A product term of input $(\beta, c, i', j')$ and flipped tap $(o, c', a, d)$ has exponent
$(\beta c_i c_o + o c_i + c_i - 1 + c - c')\,hw + (i'+a)\,w + (j'+d)$. (i) The column offset $j'+d \le w + k_w - 2$
never reaches $w + k_w - 1 + j$, so no term carries from the previous row onto an output position, and the row offset
is at most $(h + k_h - 2)w + w + k_w - 2$, so a carry from the previous block lands below $(k_h-1)w + k_w - 1$, the
smallest output offset. (ii) Within the target block $\beta c_i c_o + o c_i + c_i - 1$ only $c = c'$ contributes;
$i' + a = k_h - 1 + i$ and $j' + d = k_w - 1 + j$ select exactly the taps of the valid convolution after the flip.
(iii) Exponents reach $N$ only for $\beta = b-1$, $o = c_o - 1$, $c > c'$, i.e.\ blocks $b c_i c_o + (0 \ldots c_i - 2)$;
negacyclic reduction maps them to blocks $0 \ldots c_i - 2$ (with a sign), and output blocks are $\equiv c_i - 1
\pmod{c_i}$, so no wrapped term reaches an output coefficient. \hfill$\square$

The remaining coefficients hold partial sums of other channel pairings; they are masked like the outputs and never
decrypted. The conv-triple test (\texttt{cheetah\_conv\_triple\_test}) checks every output exactly against a plaintext
convolution, for the 23 CIFAR-10 layer shapes, the ImageNet layers and randomized shapes.

\paragraph{Tiling.} A layer needs $T = \lceil B/b\rceil\lceil (H{-}k_h{+}1)/y_h\rceil\lceil (W{-}k_w{+}1)/y_w\rceil$
tiles, $G_{\mathrm{in}} = \lceil C/c_i\rceil$ input ciphertexts and $G_{\mathrm{out}} = \lceil M/c_o\rceil$ output
ciphertexts per tile. The tiling minimizes the bytes on the wire (below),
\[
  \min_{b, h, w, c_o}\; T\cdot\bigl(G_{\mathrm{in}}\cdot \mathit{in} + G_{\mathrm{out}}\cdot \mathit{out}(b\,c_o\,y_h y_w)\bigr),
  \qquad c_i = \min\bigl(\lfloor N/(b c_o h w)\rfloor, C\bigr),
\]
where $\mathit{in}$ and $\mathit{out}(u)$ are the sizes of an input ciphertext and of an output ciphertext with $u$
used coefficients. CHEETAH's \texttt{HomConv2DSS} fixes $c_o = 1$: one output ciphertext per filter and tile. On
ImageNet's $7\times 7$ stage, such a ciphertext carries 49 useful coefficients out of 4096; the packed layout
puts up to $\lfloor N/(b\,hw)\rfloor$ filters and several images into one output ciphertext.

\begin{figure}[t]
\centering\small
\fbox{\begin{minipage}{0.95\linewidth}
\textbf{Protocol $\Pi^{\mathrm{AB2}}_{\mathrm{conv}}$} (one layer; all tiles $\tau$, input groups $g$, output groups $\gamma$)\\[2pt]
\emph{Inputs:} P0: $\lambda_W$, $\lambda_x^0$, P1's public key $\mathit{pk}_1$. P1: $\lambda_x^1$, secret key $s_1$.
\emph{Outputs:} P0: $R$; P1: $\Gamma - R$ with $\Gamma = \lambda_W \star (\lambda_x^0 + \lambda_x^1)$, $R$ uniform.
\begin{enumerate}\itemsep1pt
\item \textbf{P1 encrypts.} For each $(\tau, g)$: draw a seed $\varsigma$, expand $a = \mathrm{PRG}(\varsigma) \in R_q$,
  draw $e \leftarrow \chi$, set $c_0 = -a s_1 + e + \Delta\,\pi_x^{\tau,g}(\lambda_x^1)$. Send $(\varsigma, c_0)$:
  16 bytes and $N(60+49)$ bits.
\item \textbf{P0 evaluates.} Expand $a$ from $\varsigma$ directly in NTT form; add its own share,
  $c_0 \mathrel{+}= \Delta\,\pi_x^{\tau,g}(\lambda_x^0)$. For each $\gamma$:
  $\mathit{ct}_{\tau,\gamma} = \sum_g \pi_W^{\gamma,g}(\lambda_W)\cdot(c_0, a)_{\tau,g}$ (plaintext--ciphertext
  products, accumulated lazily in 128 bits).
\item \textbf{P0 re-randomizes and masks.} Flood: add a polynomial with uniform 64-bit coefficients to $c_0$
  (CHEETAH's $\min(64, \lceil\log q\rceil - \lceil\log t\rceil - 2)$ bits), switch to modulus $q_0$, add a fresh
  $\mathrm{Enc}_{\mathit{pk}_1}(0)$. Draw $R_{\tau,\gamma} \leftarrow R_t$
  uniformly (all $N$ coefficients) and subtract $\Delta R_{\tau,\gamma}$. P0's output share of output $(\beta,o,i,j)$
  is $R_{\tau,\gamma}[e(\beta,o,i,j)]$.
\item \textbf{P0 sends only what decryption reads.} Truncate for decryption as CHEETAH does: at $q_0$ (60 bits) drop
  the low 26 bits of $c_0$ and the low 14 bits of $c_1$, the widths below which bits only carry noise. Send $c_1$ (all $N$ coefficients, 46 bits each) and $c_0$ only at the used positions $E_{\tau}$
  (34 bits each).
\item \textbf{P1 decrypts} the coefficients $E_\tau$: $\lfloor (c_0 + c_1 s_1)_e \cdot t/q_0 \rceil$ for $e \in E_\tau$,
  its share of the outputs.
\end{enumerate}
\end{minipage}}
\caption{The AB2 convolution-triple protocol (\texttt{PackedConv2D}). AB runs it twice, once in each direction,
with P0's and P1's roles swapped for the second cross term.}
\label{fig:convproto}
\end{figure}

\paragraph{Security.} \textbf{Theorem 1.} $\Pi^{\mathrm{AB2}}_{\mathrm{conv}}$ securely computes the functionality that
outputs a uniformly random $R$ to P0 and $\Gamma - R$ to P1 against a semi-honest adversary corrupting one party,
assuming decision-RLWE for $(N, q, \chi)$, the seed expansion modeled as a random oracle, and the statistical
circuit privacy of CHEETAH's flooding with its parameters.

\emph{Proof sketch.} \emph{Corrupted P0.} Its view is the pairs $(\varsigma, c_0)$. With $a = \mathrm{PRG}(\varsigma)$
modeled as a random oracle, $a$ is uniform and public, which is exactly the RLWE setting (SEAL's seeded
ciphertexts rely on the same argument). A hybrid over the ciphertexts replaces each $c_0 = -a s_1 + e + \Delta m$ by
a uniform element; the simulator sends uniform $(\varsigma, c_0)$. \emph{Corrupted P1.} Its output
$\Gamma - R$ is uniform because $R$ is. Before truncation, the ciphertext P1 receives is, by the flooding and the fresh
encryption of zero, statistically close to a fresh $\mathrm{Enc}_{\mathit{pk}_1}(\Pi)$ of the plaintext $\Pi = P - R$,
where $P$ is the product polynomial. Since $R$ covers all $N$ coefficients, $\Pi$ is uniform in $R_t$: its used
coefficients equal P1's output and the others are independent and uniform. The simulator therefore samples
$\Pi$ uniform subject to P1's output on $E_\tau$, encrypts it with flooding, and applies the same deterministic
truncation and selection of $c_0$ coefficients. The changes of this work leave this argument intact: the
layout only changes where values sit (the tiling depends on public shapes only), and the exact wire format and the
partial $c_0$ are deterministic functions of the message CHEETAH would send, so they cannot reveal more. The
bits that are not sent are either zero after truncation (checked at run time) or coefficients that decryption
never reads. Finally, the evaluator's work does not depend on the weights' values (\cref{sec:zero}), so its running
time reveals nothing about zero weights, which matters when the weights are known in preprocessing.
\hfill$\square$

\paragraph{Communication.} An input ciphertext costs $\mathit{in} = 16 + N(60+49)/8 = 55{,}824$ bytes, against
62.4\,KB as a zstd-compressed SEAL seeded ciphertext; an output ciphertext costs
$\mathit{out}(u) = \lceil 34u/8\rceil + 46N/8 = 23{,}552 + 4.25\,u$ bytes. For each layer, the tiling minimizes
$T(G_{\mathrm{in}}\,\mathit{in} + G_{\mathrm{out}}\,\mathit{out})$; since an input costs about twice an output, it
favors fewer inputs. Two effects add up on ImageNet (\cref{tab:heformal}). Packing several filters and images into
an output ciphertext removes most of the output ciphertexts of the late stages, and the exact wire format then cuts
the remaining bytes. Together with AB2, the conv triples of a ResNet50 inference need 479\,MiB, about a quarter
of the roughly 2,000\,MiB of AB with CHEETAH's layout (s0 with \flag{A_KNOWN=0}: 2,115\,MiB of triples in total).

\begin{table}[t]
\centering\small
\begin{tabular}{@{}lrr@{}}
\toprule
ImageNet ResNet50, one image & triples, P0 sent + received (MiB) & factor\\
\midrule
CHEETAH \texttt{HomConv2DSS} ($c_o = 1$), SEAL serialization (s0) & 1,112 & 1.00\\
packed layout (s1) & 698 & 0.63\\
+ exact wire format, byte-optimal tiling (s9, s10) & 587 & 0.53\\
\midrule
conv triples only, final: AB2 (\flag{A_KNOWN=1}) / AB (\flag{A_KNOWN=0}) & 479 / 958 & 1 / 2\\
\bottomrule
\end{tabular}
\caption{Where the conv-triple traffic goes (UC2, 32-bit ring). The first three rows include the Boolean and FC material (108\,MiB,
unchanged); the last row compares the two triple types with the same protocol.}
\label{tab:heformal}
\end{table}

\paragraph{What is left.} A model of the layout search and the wire format (\texttt{docs/paper/he\_model.py}) reproduces
the measured 244.1\,MiB of input and 234.9\,MiB of output ciphertexts of the ImageNet convolutions. Dense packing in
the same format would need 234.7\,MiB, so about half of the traffic is unused coefficients. At $56\times56$ a channel
fills 3,136 of 4,096 coefficients (77\%). At $14\times14$ and $7\times7$ the layout has to trade inputs against
outputs: a polynomial holds $c_i$ input channels and $c_o$ output-channel blocks, so inputs use about $1/c_o$ and
outputs about $1/c_i$ of their coefficients (10--48\% and 4--24\% in the chosen layouts). Stages~3 and~4 send 237\,MiB
where dense packing would need 67\,MiB. Three cheap levers do not help. A per-layer ring dimension $N = 8192$ (same
109-bit modulus) is never cheaper, except on the first layer at equal cost. Fewer bits of the input ciphertexts would
raise the product noise, which P1 knows and the 64-bit flooding must hide: with 12 bits of headroom below $\Delta =
2^{77}$, at most about 16 of 109 bits could go, paid for with more output bits. The output format already keeps only
the bits decryption needs.

The remaining lever is \emph{output repacking}: the evaluator packs the used coefficients of several sparse output
ciphertexts into dense ones with ring automorphisms (a trace per ciphertext, or PackLWEs-style merging). That needs key
switching, and key switching needs a special prime on top of the 109-bit data modulus (60 + 49 bits: $2^{32}$ plaintexts
and 64-bit flooding). At $N = 4096$ the 128-bit security bound is already exhausted by those 109 bits, so repacking
forces $N = 8192$ (with a 60-bit special prime, 169 of at most 218 bits). With the layout free to use dense inputs
($c_o = 1$) and outputs packed after the evaluation, the model gives 134\,MiB of inputs and 108\,MiB of outputs,
243 instead of 479\,MiB per HE product, plus 4.7\,MiB of Galois keys that P1 sends once (13 keys, SEAL serialization).
Keeping today's layouts and repacking only the outputs would give just 420\,MiB. Over the whole preprocessing
(\texttt{repack\_estimate.py}, measured traffic of the other components) this removes 22--38\% of the traffic in UC1,
whose AB triples need both products (1,230--2,150 to 766--1,686\,MiB), and 14--29\% in UC2 (790--1,710 to
558--1,478\,MiB), or 19--32\% and 12--23\% counting the online phase as well. UC3 has no linear-layer triples and
gains nothing. The runtime cost is in the key switches. SEAL takes 583\,\textmu s for one Galois automorphism at
$N = 8192$ on the Zen~4 hosts (1,308\,\textmu s on Zen~3), single-threaded. Calibrated against the measured conv
phase, the evaluation itself stays about the same (0.64 against 0.71\,s on Zen~4), and the repacking adds 1.1\,s with a
merge tree ($7.2\cdot10^4$ key switches) or 8.2\,s with traces ($5.4\cdot10^5$) on 32 threads; on Zen~3, 1.7 or
12.4\,s. So repacking would make an inference's preprocessing about 1.0\,s slower on Zen~4 (1.8\,s on Zen~3) to save
236\,MiB per HE product. That breaks even at about 1.9\,Gbit/s (Zen~3: 1.1\,Gbit/s): it pays off on WAN links and
loses on the 25\,Gbit/s testbed.

\paragraph{Output repacking, implemented.} \flag{CHEETAH_CONV_REPACK=1} (round~4 of \cref{sec:triad}) packs with
PackLWEs' merge tree. The input polynomial holds $C$ channel slots per position, channels fastest ($C$ a power of
two), and a weight polynomial one filter, channel $c$ at $X^{C \cdot p - c}$ for kernel offset $p$; the product then
holds each output at a multiple of $C$, and every other coefficient mixes different channels, also where it wraps
around. The evaluator merges the products of $C$ filters of a tile level by level, $m = C/2, \dots, 1$: with the
automorphism $\sigma: X \mapsto X^{N/m + 1}$, which fixes $X^{2mk}$ and negates $X^{m(2k+1)}$, the node
$A + X^m B + \sigma(A - X^m B)$ doubles the used coefficients of $A$ and $B$ and cancels their other coefficients at
the multiples of $m$; after the last level filter $r$ sits at residue $r$. The products are multiplied by
$C^{-1} \bmod q$ first ($q$ is odd), which the $\log_2 C$ doublings cancel exactly, so only the key-switching noise of
the $C - 1$ automorphisms adds up, far below the flooding, which follows the packing (as does the mask, over all
coefficients). P1 sends its Galois keys once (both parties for AB triples, 1.8\,MB each). Measured on the 53
ImageNet convolutions (32 threads, all triples exact): the byte-optimal tilings take 148{,}622 automorphisms for
264\,MiB and 6.3\s of HE; pricing each automorphism at 2\,KB in the tiling search gives 26{,}000 for 285\,MiB in 2.1\s,
against 479\,MiB in 0.71\s without repacking. The model's 243\,MiB assumed dense outputs; with $C$ a power of two,
the $49 \cdot 2^k$ positions of the feature maps fill only 77\% of a polynomial. End to end, preprocessing traffic
falls by 18--31\% in UC1 and 11--24\% in UC2, for 1.3--1.7\s more preprocessing on Zen~4 (break-even about
1.2\,Gbit/s for UC2, 2.1\,Gbit/s for UC1); the online phase is unchanged.

\paragraph{Ring dimension.} Without repacking, the packed evaluator can also run in an $N = 8192$ ring with the same
109-bit modulus (\flag{CHEETAH_CONV_POLY_N=8192}): twice the slots per ciphertext, so fewer and larger ciphertexts.
The responses do not shrink with it: an output element's share of a sparse output ciphertext costs about $c_i$ times
the coefficient width, and the tiling's $c_i$ grows with $\sqrt{N}$, so an ImageNet HE product sends 629 instead of
479\,MiB ($1.31\times$). We measured it on flare / polynize in the LAN and with a WAN shaped on both hosts
(\texttt{shape\_network\_alt.sh}, 20\ms per direction and 200\,Mbit/s on every interface, reset afterwards;
\cref{tab:polyn}). $N = 8192$ loses in both: by 0.07--0.12\s in the LAN, and by 3.1--6.4\s in the WAN, where the conv
phase is bandwidth-bound and takes $1.24\times$ as long. The online phase does not depend on $N$; in the WAN it is
round-bound (1,714--1,818 rounds of about 19\ms). $N = 4096$ stays the default.

\begin{table}[t]
\centering\small
\begin{tabular}{@{}llrrrrr@{}}
\toprule
build & $N$ & LAN pre s & WAN pre s & WAN conv s & WAN online s & conv MiB\\
\midrule
UC1 A2bits RCA & 4096 & 3.20 & 36.98 & 26.74 & 32.30 & 958\\
 & 8192 & 3.28 & 43.33 & 33.25 & 32.40 & 1{,}258\\
UC2 A2bits RCA & 4096 & 3.17 & 30.39 & 16.87 & 32.12 & 479\\
 & 8192 & 3.29 & 33.52 & 20.91 & 32.19 & 629\\
UC2 reshared RCA & 4096 & 2.05 & 25.02 & 15.79 & 33.54 & 479\\
 & 8192 & 2.12 & 28.99 & 19.53 & 33.52 & 629\\
\bottomrule
\end{tabular}
\caption{Ring dimension of the packed conv triples (ImageNet, dummy weights, flare / polynize; WAN: 20\ms per
direction, 200\,Mbit/s; median of 3 WAN runs, one LAN run; conv traffic sent and received by P0).}
\label{tab:polyn}
\end{table}
